\chapter{Computation for the second Hatcher-Wagoner invariant}
\label{chap:computation-for-the-second-hatcher-wagoner-invariant}

In this chapter, first we recall the definition for the second Hatcher-Wagoner invariant, namely, $\Theta: \ker\Sigma\to \Wh_1(\pi_1M, \mathbb{Z}_2\times \pi_2M)/\chi(K_3\mathbb{Z}[\pi_1 M])$. For simplicity, we restrict to the case when the Cerf diagram contains only a single eye of (2,3)-handle pair. For general definitions and well-definedness, see \cite{AST_1973__6__1_0} and \cite{singh2022pseudoisotopiesdiffeomorphisms4manifolds}.

Recall that $\Wh_1(\pi_1M, \mathbb{Z}_2\times \pi_2 M)=(\mathbb{Z}_2\times \pi_2 M)[\pi_1 M]/(\beta\cdot [1], \alpha\cdot [\sigma]-\alpha^\tau\cdot [\tau\sigma\tau^{-1}], \alpha,\beta\in \mathbb{Z}_2\times \pi_2 M, \tau,\sigma\in \pi_1 M)$ (see \cite{singh2022pseudoisotopiesdiffeomorphisms4manifolds} for details). 

To define $\Theta$, we need the following data: For $\{f_t,t\in I\}\in \pi_1(\mathcal{F},\mathcal{E})$ with Cerf diagram only a single eye of (2,3)-handle pair, suppose that the eye appears just before $t=\epsilon$ at critical value $\frac{1}{2}$, and ends right after $t=1-\epsilon$ at the same critical value, then consider $V:=\bigcup_{t\in [\epsilon, 1-\epsilon]} f_t^{-1}(1/2)=M'\times[\epsilon, 1-\epsilon]$ where $M'=M\# S^2\times S^2$. For each $f_t, t\in [\epsilon,1-\epsilon]$, we call the belt 2-sphere of the 2-handle $B_t$, the attaching 2-sphere of the 3-handle $A_t$. Let $A:=\bigcup_{t\in[\epsilon, 1-\epsilon]} A_t\cong S^2\times I\subset V$, $B:=\bigcup_{t\in[\epsilon, 1-\epsilon]} B_t\cong S^2\times I\subset V$. Without loss of generality, we may assume that all intersections among $A,B, V$ are transverse. Then we know that $A\cap B=I\sqcup_{i=1,...,k} S^1$, we denote the $i$-th $S^1$ as $S_i^1$. Now, choose a base point $*\in X\times 0\subset X\times I$, a fixed $*'\in I\subset A\cap B$ and a path $\delta_0\subset X\times I:$ $*\to *'$. For each $i$, choose a point $p_i\in S_i^1$, and two paths: $\delta^A_i\subset A$: $*'\to p_i$,
$\delta^B_i\subset B: *'\to p_i$. \emph{Denote} $\gamma_i:=\delta^B_i*(\delta^A_i)^{-1}\in \pi_1(X\times I, *')$ and $\gamma_i^*:=\gamma_i^{\delta_0}\in \pi_1(X\times I,*)$ by pulling $\gamma_i$ from the base point $*'$ back to the base point $*$ along $\delta_0$. Also note that both $A$ and $B$ are simply connected, then choose $D_i^A\looparrowright A$ such that $\partial D^A_i=S^1_i$, $D_i^B\looparrowright B$ such that $\partial D^B_i=S_i^1$, so \emph{denote} $\beta_i:=[D^A_i\cup D^B_i]\in \pi_2(X\times I, p_i)$  and $\beta_i^*:=\beta_i^{\delta_0*\delta_i^B}\in \pi_2(X\times I,*)$ by pulling $\beta_i$ from base point $p_1$ back to base point $*$ along $\delta_0*\delta_i^B$. 

To define the $\mathbb{Z}_2$ component, we consider two framings of the same $\mathbb{R}^2$-vector bundle on $S_i^1$. Consider $\nu (S_i^1, B)$(the normal bundle of $S_i^1$ in $B$), since $A$ intersects $B$ transversely in $V$, then $E=\nu (S_i^1, B)=\nu (A,V)|_{S_i^1}$ is a 2-dimensional trivial vector bundle over $S_i^1$. We define two framings on $E$: $e_{i,B}$ is the framing naturally induced by $D^B_i\looparrowright B$, $e_{i,A}$ is induced by the attaching data for the 3-handle, namely, $A_t\times D^2\hookrightarrow M'=M'_t$. Then $e_{i,A}-e_{i,B}\in \pi_1(SO_2)=\mathbb{Z}$, we define $s_i:=e_{i,A}-e_{i,B} \text{ mod }2$. Thus for every $i$, we have defined $s_i\in \mathbb{Z}_2$.

Using all those notations above, we can define the second Hatcher-Wagoner invariant $\Theta$ for $F:=\{f_t,t\in I\}\in \pi_1(\mathcal{F},\mathcal{E})$:
\begin{definition}
    For $F:=\{f_t,t\in I\}\in \pi_1(\mathcal{F},\mathcal{E})$ with Cerf diagram only a single eye of (2,3)-handle pair, define $\Theta(F):=\sum_{i=1,...,k} (s_i, \beta_i^*)\cdot [\gamma_i^*]$. 
\end{definition}

\begin{example} \label{eg:calculation-for-delta_k}
    Now we calculate the second Hatcher-Wagoner invariant for implanted barbell diffeomorphism $\delta_k$: 
    
    Using results in last chapter, by projecting the movements to a single $S^1\times D^3\# S^1\times S^3$, and then using the dotted version and \emph{one-parameter version of dotted and 0-framed replacement} in $S^1\times D^3$, we use $\beta_t$ and $\gamma_t^\bullet$ in $S^1\times D^3$ to represent $A_t, B_t$ in $S^1\times D^3\# S^1\times S^3$. We see that the attaching sphere $A_t=\beta_t$ intersects $B_t=\gamma_t^\bullet\cup D^2$ at a single point for $t\in[0,2]$ and $\bigcup _{t\in[2,3]} A_t\cap B_t=\beta_2^\bullet \cap \gamma_2^\bullet=I\sqcup S^1$ with $I$ the blue arc in \Cref{fig:mid-ball-transformation}, and $S^1$ the green one in \Cref{fig:mid-ball-transformation}. Thus from \Cref{fig:mid-ball-transformation} we see $\gamma^*=\gamma_1^*=(k-1)^*=k-1\in \pi_1(S^1\times D^3, *)=\mathbb{Z}$. 

    To calculate $\beta^*$, we need to find $D^B\looparrowright \gamma_2^\bullet \times I $ and $D^A\looparrowright \beta_2^\bullet $ which bound $\partial D^A=\partial D^B=S^1$. By projecting to a single $S^1\times D^3$ for $t\in[2,3]$, we know that $\gamma_t^\bullet=\gamma_2^\bullet$ is the standard green one in \Cref{fig:mid-ball-transformation} which won't change, so it's enough to find $D^B\looparrowright \gamma_2^\bullet $. For $D^B$, we have an immediate choice in \Cref{fig:mid-ball-transformation}, which is just the small disk in green $\gamma_2^\bullet$ which bounds the green $S^1$. For $D^A$, we pull the green $S^1$, which is the meridian of the brown tube, along the black tube back to the meridian of brown $S^1\subset $ the blue $D^3$, which is just a parallel of red $S^1\subset$ the blue $D^3$, then capping with the red $D^2$ which is the red long tube connected with $R_0$. In all, $\beta^*=\beta_1^*=(R_0^\gamma)^*\in \pi_2(S^1\times D^3, *)=0$, here $\gamma$ is the arc in the definition of barbell, which is the arc connecting $S$ with $R_0$.

    Now we calculate $s=s_1$. For $s_1^A=\text{attaching framing of }\nu (A, V)|_{t\in [2,3], S^1\subset A}$, from the dotted version we know that $\{A_t|t\in [2,3]\}\subset \beta_2^\bullet$ is 0-framed, thus $s_1^A$ has a chapter which is the normal bundle of $\beta_2^\bullet$ in $S^1\times D^3$, which is exactly a normal chapter of $\nu (S^1, \gamma_2^\bullet)\subset \nu (S^1, \gamma_2^\bullet \times I=B|_{t\in [2,3]})$ when restricted to the green $S^1$. Thus $s^A=s^B$, $s=0$.

    Thus we obtain that $\Theta(\delta_k)=(0, (R_0^\gamma)^*)[k-1]=(0,0)[k-1]=0\in \Wh_1(\mathbb{Z}, \mathbb{Z}_2\times 0)$.    
\end{example}

The calculation can be generalized similarly to any half-unknotted implanted barbell $\beta=(R_0, S, \gamma)$ with $S$ unknotted. After doing an isotopy of $M$ we can put $S=\partial \beta_0^\bullet$ into the standard position and use strategies in \Cref{chap:changing-from-12-handle-pair-to-23-handle-pair} to find a Cerf diagram containing a single eye of (2,3)-handle pair with $\gamma_t^\bullet, \beta_t, t\in [0,3]$. Note that for $t\in [0,2]$, $\beta_t$ intersects $\gamma_t^\bullet$ at a single point, so as the example suggests, the second Hatcher-Wagoner invariant can be calculated just using $\beta_2^\bullet$ and $\gamma_2^\bullet=\gamma_0^\bullet$. Assume we have obtained $\beta_2^\bullet\subset M$, $\beta_2^\bullet \cap \gamma_2^\bullet=I\sqcup _{i=1,...,k} S^1$, then the same argument of the example works in general case, i.e. using notations above, we have
$$\Theta(F)=\sum_{i=1,...,k}(s_i, \beta_i^*)\cdot [\gamma_i^*]$$
where for each $i$, $\gamma_i=\delta^B_i*(\delta^A_i)^{-1}$ with $\delta^A_i\looparrowright \beta_2^\bullet$, $\delta^B_i\looparrowright \gamma_2^\bullet$ connecting $*'$ and $p_i\in S^1_i$, $\beta_i=D^A_i\cup D^B_i$ with $D_i^A\looparrowright \beta_2^\bullet$ and $D^B_i\looparrowright \gamma_2^\bullet$, $s_i=e_{i,A}-e_{i,B} \text{ mod }2$ with $E=\nu (S_i^1, \gamma_2^\bullet \times I)=\nu (\beta_2^\bullet , M\times I)|_{S_i^1}$. $\nu (S_i^1, \gamma_2^\bullet \times I)=\nu (S^1_i, \gamma_2^\bullet)\oplus E'=\nu (S_i^1, D_i^B)\oplus E'$, $\nu (\beta_2^\bullet, M\times I)|_{S_i^1}=\nu (\beta_2^\bullet, M)\oplus E''$ where $E',E''$ are two 1-dimensional trivial bundle over $S_i^1$. $e_{i,B}$ is the framing induced by $\nu (S_i^1, \gamma_2^\bullet)$ and $e_{i,A}$ is the framing induced by $\nu (\beta_2^\bullet,M)|_{S_i^1}$. But since $\beta_2^\bullet$ intersects $\gamma_2^\bullet$ transversely at $S_i^1$ in $M$, so $\nu (S_i^1, \gamma_2^\bullet)=\nu (\beta_2^\bullet,M)|_{S_i^1}$, thus, $s_i=e_{i,A}-e_{i,B}=0$.

To calculate $\beta_i^*$ and $\gamma_i^*$, we need a smarter change: It can be very hard to visualize $\beta_2^\bullet$, especially when $R_0$ goes through $\beta_0^\bullet$ in a strange way ($R_0\cap \beta_0^\bullet$ may be knotted in $D^3=\beta_0^\bullet$). But we know that there is an isotopy $\phi_2: M\to M$ which is isotopic to identity such that $\phi_2(\beta_1^\bullet=\beta_0^\bullet)=\beta_2^\bullet$, $\phi_2(\gamma_1^\bullet)=\gamma_2^\bullet$. This means we can pullback every data we need to $\beta_1^\bullet$ and $\gamma_1^\bullet$. Then it will be easier to visualize in the given implanted barbell $\beta$. 

Note that $\beta_1^\bullet=\beta_0^\bullet$ is the given unknotted ball which bounds $S$, i.e. $S=\partial\beta_1^\bullet$. Also recall that $\gamma_1^\bullet=\gamma_0^\bullet \text{ tube}_\gamma \text{ }R_0$. Without loss of generality, we may assume that $\beta_0^\bullet \cap R_0=\sqcup_{i=1,...,k} S_i^1$, $\text{int}(\beta_0^\bullet)\cap \gamma=\emptyset  $ (by finger-pushing $R_0$ along $\gamma$ without reaching $S$, such that the arc $\gamma$ of the implanted barbell $\beta$ can be made arbitrarily short). Let $*_0=S\cap \gamma$ be the base point, note that $*_0$ will not move during the whole time $t\in [0,3]$(in \Cref{fig:mid-ball-transformation}, $*_0$ is the left vertex of the blue $I$). For each $i$, find a path $\delta^B_i\subset \beta=(R_0, S,\gamma)$ which is a path from $*_0$ to $p_i\in S_i^1$. Also, for each $i$, $S_i^1$ divides $R_0$ into two embedded disk $D_i$ and $D_i'$, where $D_i'$ is the one connected to the arc $\gamma$. Let $D^B_i=D_i$. Now comes the following main theorem:

\begin{theorem} \label{thm:main-calculation}
    For a half-unknotted implanted barbell $\beta=(R_0, S, \gamma)$ with $S=\partial \beta_0^\bullet$ where $\beta_0^\bullet: D^3\hookrightarrow M$, by finger-pushing $R_0$ along the arc, we can make $\gamma$ short enough such that $\text{int}(\beta_0^\bullet)\cap \gamma =\emptyset$. Now suppose that $\beta_0^\bullet \cap R_0=\bigsqcup _{i=1}^k S_i^1$. Choose $p_i\in S_i^1$ and let $*_0=\gamma\cap S$ be the base point. For each $i$, find a path $\delta^B_i\subset \beta=(R_0, S,\gamma)$ which is a path from $*_0$ to $p_i\in S_i^1$. Also, for each $i$, $S_i^1$ divides $R_0$ into two embedded disks $D_i$ and $D_i'$, where $D_i'$ is the one connected to the arc $\gamma$. Let $D^B_i=D_i$. Then the $f_\beta\in \pi_0\mathcal{P}$ we constructed in \Cref{cor:construction-of-fbeta} which results in the implanted barbell diffeomorphism w.r.t. $\beta$ satisfies:
    
    $$\Theta(f_\beta)=\sum_{i=1,...,k} (0, [D^B_i]^{\delta_i^B})\cdot  [\delta_i^B]$$
    Here we identify $\pi_i(M,*_0)$ with $\pi_i(M, \beta_0^\bullet)$ so that $[D_i]\in \pi_2(M,\beta_0^\bullet)$ and $\delta_i^B\in \pi_1(M,\beta_0^\bullet)$. 
\end{theorem}

\begin{proof}
    In the above discussion we already showed that $s_i=0\in \mathbb{Z}_2$ for all $i$, we only need to show that $[D_i]^{\delta_i^B}=\beta_i^*$ and $[\delta_i^B]=\gamma_i^*$. But since we choose $*_0\in M$ which won't move during the whole procedure, then $*=(*_0,0)\in M\times 0$, $*'=(*_0,1/2)\in (M')_t=M\# S^2\times S^2$. So for the calculation, we can choose $\beta_i^*$ and $\gamma_i^*$ based at $*_0\in M$. Then by definition and by changing data to $\gamma_1^\bullet=\gamma_0^\bullet \text{ tube}_\gamma \text{ }R_0$ and $\beta_1^\bullet=\beta_0^\bullet$, we pick $\delta_i^A\subset \beta_1^\bullet \text{ connecting }$ $*_0$ with $p_i$, pick $D_i^A\looparrowright \beta_1^\bullet$ with $\partial D_i^A=S_i^1$. Then by definition $\beta_i^*=[\phi_2((D_i^A\cup D_i^B)^{\delta_i^B})]=[(D_i^A\cup D_i^B)^{\delta_i^B}]\in \pi_2(M, *_0)$, since by \Cref{lem:one-parameter-replacement}, $ \phi_t: M\to M,t\in[1,2]$ is a one-parameter family of isotopies of $M$ with $\phi_1=\id$ and $\phi_t(*_0)=*_0,\forall t\in [1,2]$. But $\beta_0^\bullet$ is an embedded contractible $D^3$, so $\beta_i^*=[D_i^B]^{\delta_i^B}$. Also for the same reason, $\gamma_i^*=\delta_i^B*(\delta_i^A)^{-1}\in \pi_1(M, *_0)$ which is $\delta_i^B\in \pi_1(M, \beta_0^\bullet)$. 
\end{proof}

\begin{example}
    Using \Cref{fig:dotted-version,fig:move-disk} (where we draw $\gamma_1^\bullet$ and $\beta_1^\bullet=\beta_0^\bullet$ explicitly when $\beta=\delta_k$ in $S^1\times D^3$) and results in \Cref{thm:main-calculation}, one can easily see that for the implanted barbell $\delta_k$ in $S^1\times D^3$, $\Theta(f_{\delta_k})=(0,0)\cdot [k-1]=0$, which coincides with \Cref{eg:calculation-for-delta_k}.
\end{example}
